\documentclass[a4paper,10pt]{article}

\usepackage[T1]{fontenc}
\usepackage[utf8]{inputenc}
\usepackage{array}
\usepackage{listings}
\usepackage{url}
\Urlmuskip=0mu plus 1mu

% A compact article layout using only packages present in a minimal TeX setup.
\setlength{\textwidth}{6.83in}
\setlength{\oddsidemargin}{-0.28in}
\setlength{\evensidemargin}{-0.28in}
\setlength{\textheight}{9.9in}
\setlength{\topmargin}{-0.35in}
\setlength{\headheight}{0pt}
\setlength{\headsep}{0pt}
\setlength{\parindent}{1em}
\setlength{\parskip}{0.15em}
\setlength{\intextsep}{0.6em}
\setlength{\textfloatsep}{0.8em}
\makeatletter
\renewcommand\section{\@startsection{section}{1}{\z@}%
  {-1.4ex plus -.3ex minus -.2ex}{0.6ex plus .2ex}{\large\bfseries}}
\renewcommand\subsection{\@startsection{subsection}{2}{\z@}%
  {-1.0ex plus -.2ex minus -.1ex}{0.35ex plus .1ex}{\normalsize\bfseries}}
\makeatother
\emergencystretch=2em
\widowpenalty=10000
\clubpenalty=10000

\lstset{
  basicstyle=\ttfamily\small,
  columns=fullflexible,
  keepspaces=true,
  showstringspaces=false,
  frame=single,
  breaklines=true,
  aboveskip=0.5em,
  belowskip=0.5em
}

\title{Making Stack Comments Checkable:\\Static Type Checking for Forth}
\author{Jaanus Pöial\\Tallinn University of Technology}
\date{}

\begin{document}
\maketitle
\vspace{-1.5em}

\begin{abstract}
Forth programmers already document words with stack-effect comments.  This
paper shows how nearly the same notation can drive a static checker.  An effect
records the required and produced stack suffix, while indexed type names retain
the identity of values moved by words such as \texttt{DUP}, \texttt{SWAP}, and
\texttt{OVER}.  Three operations form the useful core: join adjacent effects
for straight-line code, find a common contract for alternatives, and compare
one with two executions to approximate repetition.  A native Gforth prototype
loads its type hierarchy, word effects, parser behavior, and control structures
from text files.  It infers effects for colon definitions and reports stack or
type clashes without running the program.  The presentation emphasizes what a
Forth user can write, see, and rely on; formal proofs and deeper algebra are
left to the cited technical work.
\end{abstract}

\noindent\textbf{Keywords:} Forth; stack effects; static checking; symbolic
execution; concatenative languages.

\section{From comments to contracts}

A stack comment is one of Forth's best documentation tools:
\begin{lstlisting}
: SQUARED  DUP * ;        ( n -- n )
\end{lstlisting}
The comment states a calling convention, but a normal Forth system neither
checks it nor derives it.  A later edit can leave the comment wrong.  The
problem becomes harder around stack shuffles, branches, and loops, where stack
height alone is not enough.

The checker described here treats an effect as a small symbolic execution of a
word.  As in ordinary Forth notation, the top of stack is on the right:
\begin{lstlisting}
+       ( n n -- n )
@       ( a-addr -- x )
C@      ( c-addr -- char )
\end{lstlisting}
The effect constrains only the part of the data stack touched by the word.  An
unmentioned prefix passes through unchanged.  Consequently, the checker need
not know the complete stack at every point: it can infer which suffix must have
been present before a phrase.

This is a source-level checker, not a modified Forth virtual machine.  It reads
source text and computes effects instead of executing the words.  Its job is to
answer practical questions:
\begin{itemize}
  \item Does every consumer receive enough stack items of compatible types?
  \item Which input value reaches each output after stack shuffling?
  \item Do alternatives admit one usable stack contract?
  \item Does a repeated phrase appear stack-stable under the implemented test?
\end{itemize}

The approach does not make Forth dynamically typed at run time, nor does it
attach tags to cells.  Types exist only in the checker and may be as broad or
as detailed as the selected application profile.  The calculus follows earlier
work on algebraic and multiple stack effects, typed wildcards, must-analysis,
and executable checking tools~\cite{poial1991,poial1992,poial2002,poial2006,poial2008}.

\section{Effects that remember values}

\subsection{Types supplied by the application}

The checker loads a finite subtype relation.  A typical Forth-oriented fragment
is:
\begin{lstlisting}
type X cell x
type n number N
type char c
type flag boolean f
type addr a
type c-addr ca
type a-addr aa

rel n < X
rel addr < X
rel char < n
rel c-addr < addr
rel a-addr < c-addr
\end{lstlisting}
Here \texttt{X}, \texttt{cell}, and \texttt{x} are aliases.  A
\texttt{a-addr} is acceptable where a \texttt{c-addr}, \texttt{addr}, or
\texttt{X} is required, but the converse is not necessarily safe.  The
checker keeps the more specific type when two compatible boundaries meet.  It
reports a clash when neither type is a subtype of the other.

The hierarchy is deliberately profile-dependent.  For example, one profile
may place \texttt{flag} below \texttt{n}, while a stricter Forth-2012
profile~\cite{forth2012} keeps flags separate from numeric types.  This makes assumptions visible
instead of embedding one machine model in the analyzer.

\subsection{Indexed symbols}

Types alone lose important information.  The two occurrences produced by
\texttt{DUP} are not unrelated values, and \texttt{SWAP} does not manufacture
new values.  Indexed symbols state these correlations:
\begin{lstlisting}
DUP   ( X[1] -- X[1] X[1] )
SWAP  ( X[2] X[1] -- X[1] X[2] )
OVER  ( X[2] X[1] -- X[2] X[1] X[2] )
\end{lstlisting}
An index is local to an effect declaration.  Repeated \texttt{[1]} means the
same symbolic input has flowed to those positions.  Two plain, unindexed
\texttt{X} occurrences are fresh values; they are not silently assumed equal.
Before checking a phrase, the implementation renames local indices so effects
from different words cannot collide.

Indices are useful even on an untyped cell machine.  They let a use later in a
phrase refine an earlier value and carry that information through intervening
stack words.  They express data flow without introducing variable names into
the Forth source.

\section{The effect calculus in working terms}

Only three combining operations are needed for the supported control forms.
Table~\ref{tab:operations} gives their programmer-facing interpretation.

\begin{table}[h]
\centering
\small
\begin{tabular}{@{}>{\bfseries}p{0.18\linewidth}p{0.26\linewidth}p{0.48\linewidth}@{}}
\hline
Source form & Operation & Question answered \\
\hline
\texttt{P Q} & composition & Can the output of \texttt{P} supply the input of \texttt{Q}? \\
branches & common contract & What stack effect is valid for either alternative? \\
repetition & one-versus-two test & Is the repeated phrase stable in this local approximation? \\
\hline
\end{tabular}
\caption{The three core effect operations.}
\label{tab:operations}
\end{table}

\subsection{Sequence: match the touching boundaries}

For adjacent phrases, the checker matches the top output of the first effect
against the top input of the second, then works inward.  Compatible symbols
are replaced everywhere by one symbol with the more precise type.  Any
unmatched requirement becomes part of the inferred input; any unmatched result
becomes part of the output.

For example, integer literals, \texttt{+}, and \texttt{DUP} may be declared as:
\begin{lstlisting}
literal integer ( -- n )
+               ( n n -- n )
DUP             ( X[1] -- X[1] X[1] )
\end{lstlisting}
The phrase
\begin{lstlisting}
1 2 + DUP
\end{lstlisting}
first produces two numbers, consumes them with \texttt{+}, and then matches the
resulting \texttt{n} against \texttt{DUP}'s general \texttt{X}.  The inferred
result is:
\begin{lstlisting}
( -- n[1] n[1] )
\end{lstlisting}
Both output cells come from the one numeric result.  Substitution across the
whole working phrase is the key: refinement discovered at one boundary remains
visible on every correlated occurrence.

A mismatch is equally direct.  Given \texttt{C@ ( c-addr -- char )} and
\texttt{! ( X a-addr -- )}, the phrase \texttt{C@ !} attempts to use the
produced \texttt{char} as the address consumed by \texttt{!}.  If those types
are incomparable in the selected profile, checking stops at that boundary
rather than waiting for a bad memory access.

\subsection{Alternatives: find one honest contract}

An \texttt{IF} has more than one path, but callers need one effect for the
whole definition.  The checker therefore combines branch effects only when it
can construct a common contract.  It aligns their visible inputs and outputs,
accounts for a preserved stack prefix when depths differ, and unifies
corresponding types and identities.  The condition's \texttt{flag} consumption
is then composed with that contract.

The easy case is two branches with the same behavior:
\begin{lstlisting}
: KEEPIF  IF DUP ELSE DUP THEN ;
\end{lstlisting}
The inferred effect is:
\begin{lstlisting}
KEEPIF  ( X[1] flag -- X[1] X[1] )
\end{lstlisting}
By contrast, one branch leaving an extra item while the other removes one will
normally have no common usable boundary and is rejected.  Differences in type
can still be accepted when one type is a declared subtype of the other; the
result records the stronger requirement that is valid for both paths.

This operation is sometimes called a greatest lower bound in the effect
literature.  For day-to-day use, ``common branch contract'' is the important
idea.  It is not a union of everything either branch might do.  It keeps only a
single symbolic interface that the checker can justify for both.

\subsection{Loops: compare one pass with two}

A loop body can execute an unknown number of times.  The prototype uses a
finite, deliberately modest test: compute the effect of one pass, compute the
effect of two consecutive passes, and ask the branch-contract operation for a
common stable summary.  This catches common mistakes such as a body that grows
or shrinks the data stack on each iteration.

For example, including the terminating test in the repeated phrase makes this
loop stack-stable:
\begin{lstlisting}
: ULOOP  BEGIN DUP 0= UNTIL ;
\end{lstlisting}
With \texttt{0= ( n -- flag )} and \texttt{UNTIL ( flag -- )}, its inferred
effect is:
\begin{lstlisting}
ULOOP  ( n[1] -- n[1] )
\end{lstlisting}
The number survives each pass; the duplicated copy is consumed by the test.
The specification language can similarly describe \texttt{BEGIN--WHILE--REPEAT}
and \texttt{DO--LOOP} skeletons.

The limitation matters: agreement between one and two passes is a local
approximation, not a general fixed-point computation and not a proof for every
iteration count.  Acceptance says that the phrase passed this stack-effect
test under its primitive declarations.  It does not establish termination or
all semantic loop invariants.

\section{Turning Forth source into effects}

The algebra alone cannot read real source.  Forth words may consume the next
name, scan to a delimiter, change compilation state, or introduce dictionary
entries.  The prototype therefore separates three inputs:
\begin{enumerate}
  \item a \emph{type file} containing names, aliases, and subtype edges;
  \item a \emph{specification file} containing word effects and parsing roles;
  \item the Forth \emph{program file} to check.
\end{enumerate}
Representative specifications are:
\begin{lstlisting}
"("      parse until ")"       ( -- )
:        parse word define colon ( -- )
;        control end             ( -- )
literal integer                  ( -- n )
IF       control if              ( flag -- )
@                                  ( a-addr -- X )
\end{lstlisting}
Control syntax can also be declared as a pattern with captured source
segments, followed by an effect recipe using sequencing, alternatives, and
repetition.  This keeps parser conventions out of the calculus and permits
custom Forth-like profiles.

A colon definition is checked when it closes, and its inferred effect is added
to the analysis dictionary for later source.  Ordinary comments are consumed
but are not trusted as declarations.  Thus this source cannot force an
incorrect result merely by claiming one:
\begin{lstlisting}
: BAD  C@ ! ;   ( -- n )
\end{lstlisting}
Recognized locals syntax may provide a provisional shape, but a normal stack
comment remains documentation.

The bundled launchers run either the native Gforth checker or the Java
reference.  For example:
\begin{lstlisting}
./run-evaluator-gforth.sh real

./run-evaluator-gforth.sh \
  --types forth2012types.txt \
  --specs forth2012specs.txt \
  --prog forth2012prog.txt
\end{lstlisting}
On Windows, \texttt{run-evaluator-gforth.bat} provides the corresponding
entry point.  The native implementation itself is the single file
\texttt{gforth-evaluator.fs}.

\section{A complete small example}

The bundled \texttt{real} profile analyzes these definitions and a top-level
phrase:
\begin{lstlisting}
: PEEK2   DUP @ SWAP @ ;
: MAYSWAP IF SWAP THEN ;
: NLOOP   BEGIN DUP 0= WHILE REPEAT ;
: ULOOP   BEGIN DUP 0= UNTIL ;
: KEEPIF  IF DUP ELSE DUP THEN ;

DP DP C@ 0= KEEPIF
\end{lstlisting}
The generated log includes:
\begin{lstlisting}
PEEK2   ( a-addr[1] -- X[2] X )
MAYSWAP ( X[1] X[1] flag -- X[1] X[1] )
NLOOP   ( n[1] -- n[1] )
ULOOP   ( n[1] -- n[1] )
KEEPIF  ( X[1] flag -- X[1] X[1] )
\end{lstlisting}
Several details are worth reading as a Forth programmer.  \texttt{PEEK2}
requires an address because its input eventually reaches \texttt{@}; the type
requirement travels backward through \texttt{DUP} and \texttt{SWAP}.
\texttt{MAYSWAP} can promise one output order only under the stronger symbolic
condition shown by its repeated index.  Both loop examples preserve their
numeric input according to the one-versus-two rule.  \texttt{KEEPIF} consumes a
flag and duplicates the remaining cell regardless of the selected branch.

For the final phrase, the checker reports the per-word refined effects and
ends with:
\begin{lstlisting}
DP      ( -- a-addr[1] )
DP      ( -- a-addr )
C@      ( a-addr -- char )
0=      ( char -- flag )
KEEPIF  ( a-addr[1] flag -- a-addr[1] a-addr[1] )
< a-addr[1] a-addr[1]
\end{lstlisting}
The second \texttt{DP} result is specific enough for \texttt{C@}; \texttt{0=}
produces the flag consumed by \texttt{KEEPIF}; and the first address is the
value duplicated on exit.  This is more information than a stack-depth pass
would retain, yet the result still looks like familiar Forth notation.

\section{What acceptance does---and does not---mean}

The checker is useful as an early consistency check, especially for legacy
code whose stack comments are incomplete.  It can expose underflow-like shape
errors, incompatible address and scalar uses, branch interfaces that disagree,
and many loop bodies with changing stack shape.  External profiles also make
it possible to tighten a project gradually instead of imposing one universal
Forth type system.

Its guarantees are bounded by the model:
\begin{itemize}
  \item Primitive effects are trusted.  A wrong specification can make a wrong
        program appear valid.
  \item The current type match handles equal or directly comparable types; it
        does not search the hierarchy for an arbitrary third common subtype.
  \item Internal normalization can discard some automatically introduced
        identity correlations, reducing precision in later checks.
  \item Branch merging produces one conservative syntactic contract, not a set
        of all possible branch effects.
  \item The loop rule is the one-versus-two approximation described above.
  \item Return-stack tricks, unrestricted \texttt{EXECUTE}, non-local
        \texttt{THROW}, self-modifying definitions, and other metaprogramming
        may need special models or remain outside the supported subset.
\end{itemize}
The tool should therefore complement tests, reviews, and target-system checks,
not replace them.  It is best viewed as executable stack-effect documentation
with useful static diagnostics.

\section{Conclusion}

Forth already has the right surface notation for a lightweight static
analysis.  Making stack comments executable requires two additions: a
project-specific type relation and indices that preserve symbolic value flow.
From there, three operations cover much useful source: compose effects for a
word sequence, form a common contract for alternatives, and compare one with
two passes for repetition.

The result remains recognizably Forth.  Word effects live in text files,
custom parser and control behavior can be declared, definitions receive
inferred stack comments, and diagnostics point to incompatible boundaries.
The method does not solve arbitrary Forth metaprogramming or general loop
verification, but it offers a practical path from informal comments to
machine-checked stack discipline.

\begin{thebibliography}{9}
\small
\raggedright

\bibitem{forth2012}
Forth 200x Standardisation Committee.
\newblock \emph{Forth 2012 Standard}, draft proposed standard, 2014.
\newblock \url{https://forth-standard.org/standard/intro}.

\bibitem{poial1991}
J. Pöial.
\newblock Algebraic specifications of stack effects for Forth programs.
\newblock In \emph{1990 FORML Conference Proceedings, EuroFORML'90},
pp. 282--290, 1991.

\bibitem{poial1992}
J. Pöial.
\newblock Multiple stack-effects of Forth programs.
\newblock In \emph{1991 FORML Conference Proceedings, euroFORML'91},
pp. 400--406, 1992.

\bibitem{poial2002}
J. Pöial.
\newblock Stack effect calculus with typed wildcards, polymorphism and
inheritance.
\newblock Abstract and presentation, 18th EuroForth Conference, 2002.

\bibitem{poial2006}
J. Pöial.
\newblock Typing tools for typeless stack languages.
\newblock In \emph{Proceedings of the 22nd EuroForth Conference},
pp. 40--46, 2006.

\bibitem{poial2008}
J. Pöial.
\newblock Java framework for static analysis of Forth programs.
\newblock In \emph{Proceedings of the 24th EuroForth Conference},
pp. 20--23, 2008.

\end{thebibliography}

\end{document}
